\clearpage
\section{Tips for Correct Usage}
\label{sec:tips}

Using the Pi-Radio FR3-1ch is easy: pick your FR1 and FR3 center frequencies (and, on TX, the sideband), and it just
works --- the unit works out the LO frequency, switch states, and LO-null bias on its own. A few corner cases are
worth knowing about, though, because in them your choice of configuration makes a real difference. Each is
described below, with what to do about it and where in this datasheet it is measured.

\textbf{Before going over the air, check the regulations.} Before transmitting from an antenna, make sure your
signal meets every regulation that applies where you operate it --- including the permitted frequencies,
bandwidth, transmit power, and emissions outside your channel --- and that you hold any license it requires. Staying
within the rules is the operator's responsibility, and Pi-Radio is not responsible for any regulatory violation.

\textbf{Protect the unit from static (ESD).} The unit is static-sensitive: some of the parts inside it can be damaged by an
electrostatic discharge of under 150\,V (Table~\ref{tab:summary-absmax}), far less than a person can build up just
by walking across a room, and well below what you can feel. The damage is not always obvious: a part can keep working
but with degraded performance. We strongly recommend using the system in its box, as shipped: there is then never a
reason to touch the board itself, and the connectors are the only way a discharge can reach it. So take care each time
you connect something. Work on a grounded anti-static mat, wearing a grounded wrist strap; if you have neither, at
least touch grounded metal first. Before connecting any cable to any connector, touch the cable's outer shell to
ground, so that any charge it has picked up drains away there rather than through the unit; this matters most for
cables from an antenna or from other equipment.

\textbf{Keep the input within the linear range.} Each chain is linear only up to its input P1dB, which depends on
the (FR1, FR3) combination: Sections~\ref{sec:compression-upper} and \ref{sec:compression-lower} give it for TX,
Sections~\ref{sec:rx-compression-lowside} and \ref{sec:rx-compression-highside} for RX, and the Key Specifications
page summarizes the range. Drive the unit past it and the signal distorts. Separately, never exceed the absolute
maximum input levels on the Key Specifications page.

\textbf{The lower sideband (TX) and high-side LO (RX) mirror the spectrum.} With the TX lower sideband, or equivalently
a high-side LO on RX, the signal's spectrum comes out flipped: a feature at the top of the FR1 band lands at the
bottom of the FR3 band, and vice versa (Section~\ref{sec:hardware-primer}). Most receivers can undo this (for example
by conjugating the baseband samples), but your system has to know about it.

\textbf{When to engage the AM3153 filter.} The filter is what suppresses the unwanted sideband on TX
(Section~\ref{sec:sideband-suppression}) and the image on RX (Section~\ref{sec:rx-image-rejection}); with it
bypassed, both are poor. Its cost is some conversion gain, and on RX a higher noise figure
(Section~\ref{sec:rx-noise-figure}). Engage it whenever the unwanted sideband or image matters; bypass it only when
you need the lowest possible RX noise figure and know that nothing is transmitting in the image band.

\textbf{An interferer sits on the RX image.} The RX receives both LO\,+\,FR1 and LO\,$-$\,FR1, and one of them is
the image (Section~\ref{sec:rx-image-rejection}). If a known interferer sits on the image, switch to the other LO
side: the LO, and with it the image, moves by twice the FR1 frequency, away from the interferer.

\textbf{Image rejection is weakest at low FR1.} At low FR1 the wanted signal and its image are close together, so
the filter has a harder time separating them. If image rejection matters, use a higher FR1 for the same FR3
(Section~\ref{sec:rx-image-rejection}).

\textbf{Avoid an LO at exactly twice the FR1.} When the LO is exactly twice the FR1 frequency, a spurious mixing
product lands on the wanted IF and RX image rejection drops sharply (Section~\ref{sec:rx-image-rejection}). Avoid
that exact configuration, or choose the sideband so that the LO stays at 9\,GHz or above.

\textbf{RX noise figure is worse with the LO at 14--14.5\,GHz.} With the LO at 14.0 or 14.5\,GHz, the RX noise figure
is typically 4 to 5\,dB worse than at neighboring LO frequencies (Section~\ref{sec:rx-noise-figure}). This only
matters for weak signals near the noise floor; if that is your case, switch to the other LO side, which moves the
LO by twice the FR1 frequency.

\textbf{Weak CW signals on a 10\,MHz grid.} The RX IF output carries a weak comb of tones at multiples of 10\,MHz
(Section~\ref{sec:rx-noise-figure}). It matters only for an extremely weak CW signal that lands exactly on that grid;
if yours does, move the IF (or the LO) by a few hundred kHz.

\keepnext{5\baselineskip}% the whole paragraph moves to the next page rather than split
\textbf{Phase noise rises with the LO frequency.} The LO's phase noise grows with its frequency
(Section~\ref{sec:phase-noise}). For the highest-order modulations (256-QAM), choose the FR1 and sideband that put
the LO as low as possible, or use an external low-noise LO such as the Pi-Radio OctoLO via the ``LO In'' connector.

\textbf{LO suppression depends on temperature.} The LO-null bias is calibrated at two temperatures and interpolated at
runtime from the LMX2820's die temperature (Section~\ref{sec:lo-suppression}). LO suppression is best once the unit
has reached a steady temperature after power-up, and with the LMX die between the two calibrated temperatures. Keep
the LMX die below 100\,\textdegree C in any case.

\textbf{Stay inside the design range.} The unit is rated for FR3 up to 24\,GHz and FR1 from 0.7 to 4.0\,GHz.
Measurements beyond these limits are shown throughout this datasheet, boxed in red, and the unit still works there,
but its performance outside them is not guaranteed.
